%% trajectory_contributions.tex -- which factors co-move with a dependent
%% variable over time, and at which time points.

\begin{recipe}[label=rec:trajectory-contributions]{Trajectory-Based Signed Contribution Analysis}

\recipepurpose{For entities measured repeatedly over time --- countries,
  patients, cell lines --- find which measured factors move together with a
  chosen dependent variable, how strongly over the whole trajectory, and at
  which individual time points. Contributions are signed, so co-movement and
  divergence are distinguished rather than conflated.}

\nomaintainer

\begin{requiredinputs}
  \item Repeated measurements arranged as
        \code{(n\_factors, n\_samples, n\_timepoints)} --- indicators first,
        entities second, \emph{time last}. Every entity must be measured at the
        same time points.
  \item Which indices are factors and which is the dependent. Both are
        positions in the first axis: a dependent is simply an indicator you
        chose to explain.
  \item Regularly spaced time points, if you intend to use the velocity or
        acceleration variants. Interpolating irregular sampling is not part of
        the library.
\end{requiredinputs}

\begin{prerequisites}
  \item None beyond having your measurements in that layout.
\end{prerequisites}

\begin{workflow}
  \step{Optionally min--max each series to $[0,1]$ with
        \code{normalize\_all\_trajectories}, so entities of very different
        magnitude become comparable. It rescales each series independently and
        so preserves shape; check its \code{status} output, which reports
        per-series problems rather than raising them.}
  \step{Choose a baseline. Deviations are measured from it, so it decides what
        ``co-movement'' means: \optn{baseline\_min} for growth-like variables,
        \optn{baseline\_mean} for oscillating ones, \optn{baseline\_raw} when
        zero is itself meaningful.}
  \step{Compute contributions with \code{compute\_all\_contributions}: one
        signed number per (factor, dependent, entity) plus the per-timepoint
        series behind it. This is descriptive --- it ranks factors, it does not
        test them.}
  \step{Test a pair with \code{perform\_permutation\_test} and
        \code{compute\_p\_values}. The test replaces the dependent with another
        entity's, keeping time order intact, and so asks whether \emph{this}
        entity's pairing is unusual within your panel.}
  \step{Optionally repeat on the dynamics with
        \code{compute\_velocity\_acceleration\_contributions}, which answers
        whether the two variables change, and change their rate of change,
        together.}
\end{workflow}

\examplescript{python/examples/trajectory_contributions.py}{%
  Builds a synthetic panel in which one entity's dependent really is driven by
  one of its factors and every other entity is unrelated, then runs the whole
  chain so the planted answer can be checked against the output. The repository
  ships no time-series data, which is why the example generates its own.}

\begin{codefragment}[language=Python,caption={Fragments to edit: the layout, the baseline, the pair you test}]
import numpy as np
from tensor_omics import (compute_all_contributions,
                          compute_contributions,
                          perform_permutation_test,
                          compute_p_values)

# --- indicators first, entities second, TIME LAST ----------
traj = np.asfortranarray(panel)   # (n_factors, n_samples,
                                  #  n_timepoints)

# --- 1-based indices into the first axis -------------------
factors   = np.array([1, 2, 3], dtype=np.int32)
dependent = np.array([4], dtype=np.int32)

res = compute_all_contributions(
    traj, factors, dependent,
    "baseline_mean")        # or _min / _raw
totals = res["total_contributions"]   # (n_f, n_d, n_s)

# --- is entity 1's pairing unusual in THIS panel? ----------
obs = compute_contributions(
    traj[0, 0, :].copy(), traj[3, 0, :].copy(),
    "baseline_mean")
perm = perform_permutation_test(
    traj, 1, 4, 1,          # factor, dependent, sample
    "baseline_mean", 2000, random_seed=5)
p = compute_p_values(
    obs["local_contributions"], obs["total_contribution"],
    perm["local_contributions"], perm["total_contributions"])
\end{codefragment}

\begin{parameters}
  \param{baseline\_mode}{Where each variable's deviations are measured from,
    and therefore what a positive contribution means}{%
    \optn{baseline\_min} for monotone growth variables, so deviations are
    improvements above the historical minimum; \optn{baseline\_mean} for
    oscillating ones; \optn{baseline\_raw} only when zero is a real reference}
  \param{n\_permutations}{Size of the null distribution, and with it the
    smallest reportable p}{A p can never fall below $1/n$; use at least
    \code{1000} if you intend to report values near 0.01}
  \param{random\_seed}{Makes the permutation draw reproducible}{Always set it
    for anything you will report, otherwise the p changes between runs}
  \param{factor\_indices / dependent\_indices}{Which rows of the first axis are
    treated as explaining and explained}{1-based. Nothing distinguishes a
    ``factor'' from a ``dependent'' in the data --- the roles are yours}
\end{parameters}

\begin{expectedresults}
  On the example's planted panel, the driving factor separates cleanly from the
  rest: total contribution $12.7$ against $0.85$ and $-0.56$ for the two
  unrelated factors, with $p < 0.0005$ against $0.37$ and $0.55$. If a factor
  you know to be unrelated scores comparably to one you expect to matter, look
  first at whether the tested entity's \emph{dependent} differs from the other
  entities' --- see the pitfall below --- before concluding anything about the
  factor.
\end{expectedresults}

\begin{pitfall}[Pitfall: the test only ever looks at the upper tail]
  The p-value counts permutations whose contribution is $\geq$ the observed
  one. A strongly \emph{negative} contribution --- the two variables reliably
  diverging, which is a finding --- therefore lands near $p = 1$, not near
  $0$: a total of $-3.76$ scores $p = 0.63$. Divergence cannot be declared
  significant this way. Test the negated factor if you need it.
\end{pitfall}

\begin{pitfall}[Pitfall: a p-value of exactly zero]
  The p is the raw fraction $(\text{permuted} \geq \text{observed})/n$, with no
  correction, so when no permutation reaches the observed value it is reported
  as exactly $0$. No permutation test can support that. Report it as
  $< 1/n$ instead, and remember it is a resolution limit, not evidence: with
  $2000$ permutations nothing below $0.0005$ is measurable.
\end{pitfall}

\begin{pitfall}[Pitfall: what the permutation actually holds fixed]
  The null swaps in \emph{another entity's dependent}. So the question is
  never ``are these two variables related?'' but ``is this entity's pairing
  unusual among the entities I supplied?''. Where every entity shares the same
  temporal trend, a real and strong relationship correctly scores around
  $p = 0.5$, because a swapped-in dependent carries that trend too. The panel
  is the null; choose it deliberately.
\end{pitfall}

\begin{pitfall}[Pitfall: an entity whose dependent is unlike the rest]
  The mirror image of the previous trap, and easier to fall into. If the tested
  entity's dependent behaves differently from the others --- a strong trend
  where the rest are flat --- then \emph{every} factor in that entity is being
  compared against a null built from unlike dependents, and unrelated factors
  come out significant. In a panel built that way, a pure-noise factor reached
  $p = 0.022$ and outranked the genuine driver. Check that the dependent is
  exchangeable across entities before believing any of the p-values.
\end{pitfall}

\begin{pitfall}[Pitfall: feeding velocities back into the contribution routines]
  \code{compute\_velocity\_trajectories} returns
  \code{(n\_timepoints-1, n\_factors, n\_samples)} --- the axes reordered
  relative to its own input. Passing that to
  \code{compute\_all\_contributions}, which expects
  \code{(n\_factors, n\_samples, n\_timepoints)}, raises no error: the extents
  are simply read in the wrong roles and plausible numbers come back computed
  over the wrong axes entirely. Use
  \code{compute\_velocity\_acceleration\_contributions}, which starts from
  positions and keeps the layout, unless you are transposing deliberately.
\end{pitfall}

\begin{note}
  The methods document also describes weighted combinations of factors and of
  dependents, $F^{\ast} = \sum_j \alpha_j F_j$. There is no separate routine
  for these: build the combined series yourself and pass it as another row of
  the first axis. Group-level comparisons between sets of entities are likewise
  left to the user, on the totals this recipe produces.
\end{note}

\begin{interpretation}
  A total contribution is a \emph{descriptive} quantity: its sign says whether
  the factor and the dependent deviate from their baselines together or in
  opposition, and its magnitude how consistently. Ranking factors by it is the
  primary result, and it stands on its own without any p-value.

  The per-timepoint series is where the method earns its name. A total built
  from one enormous spike and twenty quiet points means something quite
  different from the same total spread evenly across the trajectory, and only
  the series distinguishes them --- read it before interpreting the total.

  Velocity contributions say whether the two variables \emph{change} together,
  which is a different claim from moving together: two variables can sit at
  consistently high levels without their year-to-year increments aligning at
  all. Acceleration contributions mark shared turning points. Both are
  answers to narrower questions than the ambient contribution, not refinements
  of it.

  Finally, keep the descriptive and the inferential parts apart when reporting.
  The contribution describes your data; the p-value describes how your entity
  compares with the other entities you happened to supply.
\end{interpretation}

\begin{references}
  \reference{Good (2005). \emph{Permutation, Parametric, and Bootstrap Tests
    of Hypotheses}. 3rd ed., Springer.}
  \reference{Phipson \& Smyth (2010). Permutation p-values should never be
    zero: calculating exact p-values when permutations are randomly drawn.
    \emph{Statistical Applications in Genetics and Molecular Biology}
    9(1):Article 39.}
\end{references}

\end{recipe}
